\documentclass[a4paper,12pt]{book}
%\documentclass[a4paper,12pt]{article}
\usepackage{amsthm}
\usepackage{amsmath,amssymb,dsfont}


\usepackage[utf8]{inputenc}
\usepackage[english]{babel}

\usepackage{graphicx}
\graphicspath{ {img/} }

\renewcommand\qedsymbol{\includegraphics[width=1.5cm]{occhiali}}


\newtheorem{theorem}{Theorem}[section]
\newtheorem{corollary}{Corollary}[theorem]
\newtheorem{lemma}[theorem]{Lemma}
\newtheorem{definition}{Definition}


%\title{Zorzilla's Theorems\thanks{Newtork Modeling}}
%\author{Ettore Mariotti}
%%%%%%%%%%%%


\begin{document}
%\includegraphics[width=\textwidth]{Zorzilla}
%\maketitle

\begin{titlepage}
\centering
\vspace*{\stretch{1}}
\vspace*{\stretch{1}}
{written with love by\par}
{\Large Ettore Mariotti, Andrea Pittaro, Enrico Lovisotto, Davide Peron, Umberto Martinelli\par}
\vspace{1cm}
{\LARGE{ \bfseries Zorzilla's Theorems \par}}
\vspace{2cm}
\includegraphics[width=9cm]{Zorzilla}
\vspace*{\stretch{2}}
\end{titlepage}

%%%
\chapter{Markov Chains}


\section{Asymptotic considerations}

\subsection{Steady-state probabilities}

For a regular Markov Chain, we have:

$$ \lim_{n \to \infty} P_{ij}^{(n)} = \lim_{ n \to \infty} Prob[ x_n=j | x_0 =i] = \pi_j > 0 \quad \forall i, j $$

\begin{enumerate}
\item Limit exsists (not obvious)
\item Limit is indipendent of the initial state
\item Limit is strictly positive
\end{enumerate}


%%%%%
\begin{theorem}
For a regular MC (Markov Chain) with states 0,1,...,N the limit distribution $\pi$ is the unique solution of the following system of equations:

\begin{itemize} %itemize
\item $\pi_j = \sum_{k=0}^N \pi_k P_{k j} $, \qquad for $j = 0,1, ..., N$

\item $\sum_{k=0}^N \pi_k = 1$, \qquad \qquad $\pi_k \ge 0 \quad \forall k$

\end{itemize}
\end{theorem}

---
\begin{proof} Let's start easy by proving the {\bfseries existence}:

\begin{itemize}
\item$ P_{i j}^{(n)} = \sum_{k=0}^N P_{ik}^{(n-1)} P_{k j}$ \qquad / it's the developing of $P^n = P^{n-1} P$
\item$ \sum_{k=0}^N P_{ik}^{(n)} = 1 \quad \forall n$ \qquad / normalisation condition

Now let's study what happens as $ n \to \infty $
\item $\pi_j = lim_{n \to \infty} P_{ij}^{(n)} = \sum_{k=0}^N \pi_k P_{kj}$ / This shows that the system have a solution

\end{itemize}

So the solution exsist. Now let's prove {\bfseries uniqueness}:


\begin{itemize}

\item Let $x_j$ be a solution
\item $x_j = \sum_{k=0}^N x_k P_{kj} $ / by construction
\item $x_l = \sum_{j=0}^N x_j P_{jl} =  \sum_{j=0}^N ( \sum_{k=0}^N x_k P_{kj} ) P_{jl} =  \sum_{k=0}^N x_k \sum_{j=0}^N P_{kj} P_{jl} = \sum_{k=0}^N x_k P_{kl}^{(2)}$

If we apply this trick again n times we can prove by induction  that:

\item $x_j$ also satisfy $ x_j = \sum_{k=0}^N x_k P_{kj}^n $ \quad  $\forall n $

Now, as $n \to \infty$ :

\item  $x_j = \sum_{k=0}^N x_k \pi_j = \pi_j (\sum_{k=0}^N x_k) = \pi_j  \Rightarrow $ Solution is unique

\end{itemize}

\end{proof}
%%

\subsection{Classes of states}
\begin{definition}[Accessible State]

State $j$ is {\bfseries accessible} from state $i$ ($i \rightarrow j$) if $\exists n \geq 0$ such that $P_{ij}^n > 0$
\end{definition}

\begin{definition}[Communicant States]

States $i$ and $j$ are said to {\bfseries communicate} ( $ i \leftrightarrow j$ ) if

 $i \rightarrow j$ and $j \rightarrow i$
\end{definition}
%%%%

{\bfseries Proprieties:}
\begin{enumerate}
\item Reflexivity: \quad $i \leftrightarrow i$
\item Symmetry: \quad if $i \leftrightarrow j$ then $j \leftrightarrow i$
\item Transitivity: \quad if $i \leftrightarrow j$ and $j \leftrightarrow k \Rightarrow i \leftrightarrow k$
 \end{enumerate}
Communication among states is an equivalence relation.

\begin{definition}[Periodicity]
The period of state i, $d(i)$, is the GCD (greatest common denominator) of set $S_i = \{ s>0 : P_{ii}^{(s)} >0 \}$
\end{definition}

\begin{theorem} Periodicity is a class of propriety:
\begin{itemize}
\item if $i \leftrightarrow j$, then $d(i) = d(j)$
\end{itemize}
\end{theorem}
--
\begin{proof}
given $$S_i = \{ s>0 : P_{ii}^{(s)} >0 \}$$

let $$n, m >0 : P_{ij}^{(m)} >0, P_{ij}^{(m)} > 0$$

we have that $$\forall s \in S_i : P_{ii}^{(s)} > 0$$

now, for total probability theorem$$P_{jj}^{(n+s+m)} = \sum_{h, k} P_{jh}^{(n)} P_{hk}^{(s)} P_{kj}^{(m)} \geq P_{ji}^{(n)} P_{ii}^{(s)} P_{ij}^{(m)} >0 $$

$$P_{jj}^{(n+s+m)} >0 \Rightarrow n+s+m \in S_j$$

$$P_{jj}^{(n+2s+m)} \geq P_{ji}^{(n)} P_{ii}^{(s)} P_{ii}^{(s)} P_{ij}^{(m)} >0 \Rightarrow n+2s+m \in S_j$$

Now let's call $d(j) =$ g.c.d. of $S_j$

$ \Rightarrow n+s+m$ and $n+2s+m$ are integer multiples of $d(j)$

Therefore s is an integer multple of $d(j)$ ?? bo raga non capisco dai miei appunti

$\forall s \in S_i$ is integer multiple of $d(j)$ $\Rightarrow d(j)$ is a common divisor of $S_i$ and so it divides the g.c.d. of $S$, $d(i)$

$\Rightarrow d(i)$ is an integer multiple of $d(j)$.

Doing this proof again switching role of i and j we prove that $d(j)$ is an integer multiple of $d(i)$

Therefore, $d(i) = d(j)$

\end{proof}
--

\begin{definition}[Return Time]
the return time to state $i := R_i$
\end{definition}

So the probability of ever going back to state i can be written as
$$f_{ii} = \sum_{n=1}^\infty f_{ii}(n) $$

\begin{definition}[Recurrent state]
a state is recurrent if $f_{ii} = 1$
\end{definition}

\begin{definition}[Transient state]
a state is transient if  $f_{ii} < 1$
\end{definition}

\begin{definition}
$M$ is the number of returns to state $i$
\end{definition}

$$\mathbb{E}[M | x_0 = i] = \sum_{k=1}^\infty \mathbb{P}[M\geq k | x_0 = i] = \sum_{k=1}^\infty f_{ii}^k = \begin{cases}
\frac{f_{ii}}{1-f_ii}, & \mbox{if transient} \\
\infty, & \mbox{if recurrent}
\end{cases}$$

\begin{definition}[Proper r.v]
given X a finite value r.v. it is called proper if $$\lim_{x\to \infty} \mathbb{P}[x\geq a] = 0$$
\end{definition}

\begin{definition}[Improper r.v]
given X a r.v. it is called improper if  $$\lim_{x\to \infty} \mathbb{P}[x\geq a] = p_\infty \ne 0$$
\end{definition}


\begin{theorem}
$$i  \mbox{ is recurrent} \iff \sum_{n=1}^\infty P_{ii}^{(n)} = \infty$$
\end{theorem}
--
\begin{proof}
Number of time state i has been visited: $$M = \sum_{n=1}^\infty \mathds{1}\{x_n =i\} $$
The following swap from expected value and infinite sum is allowed by Fubini's Theorem$$ \mathbb{E}[M | x_0 = i] = \mathbb{E}[\sum_{n=1}^\infty \mathds{1}\{x_n = i\} | x_0 = i] = \sum_{n=1}^\infty \mathbb{E}[\mathds{1}\{ x_n = i\} | x_0 = i] = \sum_{n=1}^\infty P_{ii}^{(n)} $$
So we have shown that $\mathbb{E}[M | x_0=i] = \sum_{n=1}^\infty P_{ii}^{(n)}$.\\
Now, remembering that$$ \mathbb{E}[M | x_0 = i] = \begin{cases}
\frac{f_{ii}}{1-f_ii}, & \mbox{if transient} \\
\infty, & \mbox{if recurrent}
\end{cases}$$
the proof is concluded.
\end{proof}
%%%%%%%%%%%%%%%%%%%%%%

\begin{theorem}
(not to confuse with theorem on periodicity)
$$\mbox{if } i\leftrightarrow j \mbox{and i is recurrent} \Rightarrow \mbox{ j is also recurrent}$$
\end{theorem}
--
\begin{proof}
$$\mbox{Since } i\leftrightarrow j : \exists m,n \mbox{ such that } P_{ij}^{(n)} \mbox{ and } P_{ji}^{(m)} > 0$$
$$\mbox{let } r>0 : P_{jj}^{(m+r+n)} = \sum_{r, k} P_{jh}^{(m)} P_{hk}^{(r)} P_{kj}^{(n)} \geq P_{ji}^{(m)}  P_{ii}^{(r)}  P_{ij}^{(n)}$$
$$\sum_{l=1}^\infty P_{jj}^{(l)} \geq \sum_{r=1}^\infty P_{jj}^{(m+r+n)} \geq \sum_{r=1}^\infty P_{ji}^{(m)}  P_{ii}^{(r)}  P_{ij}^{(n)} = P_{ji}^{(m)} P_{ij}^{(n)} \sum_{r=1}^\infty P_{ii}^{(r)}$$
but since \begin{itemize}
\item$i$ is recurrent $\Rightarrow \sum_{n=1}^\infty P_{ii}^{(n)} = \infty$
\item $P_{ij}^{(n)} \mbox{ and } P_{ji}^{(m)} > 0$
\end{itemize}
$$\sum_{l=1}^\infty P_{jj}^{(l)} = \infty \Rightarrow j \mbox{ is recurrent}$$

\end{proof}




%\begin{theorem}
%Let $f$ be a function whose derivative exists in every point, then $f$ is
%a continuous function.
%\end{theorem}

%\begin{theorem}[Pythagorean theorem]
%\label{pythagorean}
%This is a theorema about right triangles and can be summarised in the next
%equation
%\[ x^2 + y^2 = z^2 \]
%\end{theorem}

%And a consequence of theorem \ref{pythagorean} is the statement in the next
%corollary.

%\begin{corollary}
%There's no right rectangle whose sides measure 3cm, 4cm, and 6cm.
%\end{corollary}

%You can reference theorems such as \ref{pythagorean} when a label is assigned.

%\begin{lemma}
%Given two line segments whose lengths are $a$ and $b$ respectively there is a
%real number $r$ such that $b=ra$.
%\end{lemma}

%\begin{proof}
%To prove it by contradiction try and assume that the statemenet is false,
%proceed from there and at some point you will arrive to a contradiction.
%\end{proof}


\end{document}
